\documentclass[DIN, pagenumber=false, fontsize=11pt, parskip=half]{scrartcl}

\usepackage[utf8]{inputenc}
\usepackage[T1]{fontenc}
\usepackage{template}
\usepackage{amsmath}

\setcounter{secnumdepth}{1} % number sections only; subsections (Exercises) unnumbered

\titlehead{\centering\small
    Worked solutions to selected exercises from\\
    G. Chrystal, \emph{Algebra: An Elementary Text-Book}, Part I (1886)\\
\rule{0.4\linewidth}{0.4pt}}

\title{}
\author{} % Oscar A. Carrera
\date{}

\begin{document}
\maketitle

\setcounter{section}{2}
\section{Theory of Quotients---First Principles of Theory of Numbers}

\subsection{Exercises III}

\begin{problems}[from=1,step=1]
    \problem Express in its simplest form---
    \[
        \frac{x^2}{x-y} + \frac{y^2}{y-x}.
    \]
    \answer
    \begin{align*}
        \frac{x^2}{x-y} + \frac{y^2}{y-x} &= \frac{x^2(y-x)}{(x-y)(y-x)} + \frac{y^2(x-y)}{(x-y)(y-x)} && \text{($\frac{a}{b} + \frac{p}{q} = \frac{qa}{qb} + \frac{pb}{qb}$)}\\ 
                                          &= \frac{x^2(y-x) + y^2(x-y)}{(x-y)(y-x)} && \text{($\frac{a}{b} + \frac{c}{b} = \frac{a+c}{b}$)}\\ 
                                          &= \frac{x^2(y-x) + y^2(x-y)}{xy-x^2-y^2+xy} && \text{(expand denominator)}\\ 
                                          &= \frac{x^2(y-x) + y^2(x-y)}{x(y-x) + y(x-y)} && \text{(factor denominator)}\\ 
                                          &= \frac{x^2(y-x) - y^2(y-x)}{x(y-x) - y(y-x)} && \text{($x-y = -(y-x)$)}\\ 
                                          &= \frac{(x^2-y^2)(y-x)}{(x-y)(y-x)} && \text{(extract common factor $y-x$)}\\ 
                                          &= \frac{(x^2-y^2)}{(x-y)} && \text{($\frac{ac}{bc} = \frac{a}{b}$)}\\ 
                                          &= \frac{(x-y)(x+y)}{(x-y)} && \text{($a^2-b^2 = (a-b)(a+b)$)}\\ 
                                          &= x+y && \text{($\frac{ac}{bc} = \frac{a}{b}$)}
    \end{align*}

    \problem Express in its simplest form---
    \[
        \frac{a}{a-b} + \frac{b}{b-a}.
    \]
    \answer 
    \begin{align*}
        \frac{a}{a-b} + \frac{b}{b-a} &= \frac{a(b-a) + b(a-b)}{(a-b)(b-a)} && \text{($\frac{a}{b} + \frac{p}{q} = \frac{aq}{bq} + \frac{bp}{bq}$)} \\ 
                                      &= \frac{a(b-a) + b(a-b)}{ab-a^2-b^2+ab} && \text{(expand denominator)}\\ 
                                      &= \frac{a(b-a) + b(a-b)}{a(b-a) + b(a-b)} && \text{(factor denominator)} \\ 
                                      &= 1 && \text{($\frac{a}{a} = 1$)}
    \end{align*}


    \problem Simplify---
    \[
        \frac{P+Q}{P-Q} - \frac{P-Q}{P+Q}, \quad\text{where}\quad P = x+y, \quad Q = x-y.
    \]
    \answer
    \begin{align*}
        \frac{P + Q}{P - Q} - \frac{P - Q}{P + Q} &= \frac{x+y+x-y}{x+y-x+y} - \frac{x+y-x+y}{x+y+x-y} && \text{(replace variables)}\\ 
                                                  &= \frac{2x}{2y} - \frac{2y}{2x} && \text{(simplify)}\\ 
                                                  &= \frac{x}{y} - \frac{y}{x} && \text{($\frac{ab}{bc} = \frac{a}{c}$)}\\ 
                                                  &= \frac{x^2 - y^2}{xy} && \text{($\frac{a}{b} - \frac{p}{q} = \frac{aq}{bq} - \frac{pb}{bq}$)}
    \end{align*}
   
    \problem Simplify---
    \[
        \frac{1 - \dfrac{x(1-y)}{x+y}}{1 + \dfrac{1-y}{x+y}}.
    \]
    \answer 
    \begin{align*}
        1-\frac{x(1-y)}{x+y} \div 1+\frac{1-y}{x+y} &= \frac{x+y}{x+y}-\frac{x(1-y)}{x+y} \div \frac{x+y}{x+y} + \frac{1-y}{x+y}\\ 
                                                    &= \frac{x+y-x+xy}{x+y} \div \frac{x+y+1-y}{x+y}\\ 
                                                    &= \frac{xy+y}{x+y} \div \frac{x+1}{x+y}\\ 
                                                    &= \frac{xy+y}{x+y} \times \frac{x+y}{x+1}\\ 
                                                    &= \frac{xy+y}{x+1} \\ 
                                                    &= \frac{y(x+1)}{x+1}\\ 
                                                    &= y
    \end{align*}

    \problem Simplify---
    \[
        \frac{\dfrac{1}{ab} - \dfrac{1}{ac} - \dfrac{1}{bc}}{\dfrac{a^2 - (b-c)^2}{a}}.
    \]
    \answer
    \begin{align*}
        \frac{\dfrac{1}{ab} - \dfrac{1}{ac} - \dfrac{1}{bc}}{\dfrac{a^2 - (b-c)^2}{a}}
            &= \left(\frac{1}{ab} - \frac{1}{ac} - \frac{1}{bc}\right) \div \frac{a^2 - (b-c)^2}{a} && \text{($\frac{a}{b} = a \div b$)}\\ 
            &= \left(\frac{c}{abc} - \frac{b}{abc} - \frac{a}{abc}\right) \div \frac{a^2 - (b-c)^2}{a} && \text{(common denominator $abc$)}\\ 
            &= \frac{c-b-a}{abc} \div \frac{a^2 - (b-c)^2}{a} && \text{($\frac{a}{b} - \frac{c}{b} = \frac{a-c}{b}$)}\\ 
            &= \frac{c-b-a}{abc} \div \frac{\big(a-(b-c)\big)\big(a+(b-c)\big)}{a} && \text{($a^2-b^2 = (a-b)(a+b)$)}\\ 
            &= \frac{c-b-a}{abc} \div \frac{(a-b+c)(a+b-c)}{a} && \text{(remove brackets)}\\ 
            &= \frac{c-b-a}{abc} \times \frac{a}{(a-b+c)(a+b-c)} && \text{($\frac{a}{b} \div \frac{p}{q} = \frac{a}{b} \times \frac{q}{p}$)}\\ 
            &= \frac{-(a+b-c)}{abc} \times \frac{a}{(a-b+c)(a+b-c)} && \text{($c-b-a = -(a+b-c)$)}\\ 
            &= \frac{-a(a+b-c)}{abc(a-b+c)(a+b-c)} && \text{($\frac{a}{b} \times \frac{p}{q} = \frac{ap}{bq}$)}\\ 
            &= -\frac{1}{bc(a-b+c)} && \text{($\frac{ac}{bc} = \frac{a}{b}$)}
    \end{align*}

    \problem Simplify---
    \[
        \left(a - \frac{b^2}{a+b}\right) \times \left(a + \frac{b^2}{a-b}\right).
    \]
    \answer 
    \begin{align*}
        \MoveEqLeft
        \left(a - \frac{b^2}{a+b}\right) \times \left(a + \frac{b^2}{a-b}\right) \\
        &= \left(\frac{a(a+b)}{a+b} - \frac{b^2}{a+b}\right) \times \left(\frac{a(a-b)}{a-b}+\frac{b^2}{a-b}\right) && \text{($\frac{a}{b} - \frac{p}{q} = \frac{aq}{bq} - \frac{pb}{bq}$)}\\ 
        &= \left(\frac{a^2+ab-b^2}{a+b}\right) \times \left(\frac{a^2-ab+b^2}{a-b}\right) && \text{($\frac{a}{c} - \frac{b}{c} = \frac{a-b}{c}$)}\\ 
        &= \frac{(a^2+ab-b^2)(a^2-ab+b^2)}{(a+b)(a-b)} && \text{($\frac{a}{b} \times \frac{p}{q} = \frac{ap}{bq}$)}\\ 
        &= \frac{a^4-a^2b^2+2ab^3-b^4}{(a+b)(a-b)} && \text{(expand numerator)}\\ 
        &= \frac{a^4-a^2b^2+2ab^3-b^4}{a^2-b^2} && \text{(distribute denominator)}
    \end{align*}

    \problem Simplify---
    \[
        \frac{1}{x+y} + \frac{1}{x-y} - \frac{2x}{x^2+y^2}.
    \]
    \answer 
    \begin{align*}
        \MoveEqLeft
        \frac{1}{x+y} + \frac{1}{x-y} - \frac{2x}{x^2+y^2}\\ 
        &= \frac{x-y+x+y}{(x+y)(x-y)} - \frac{2x}{x^2+y^2} && \text{($\frac{a}{b} - \frac{p}{q} = \frac{aq}{bq} - \frac{pb}{bq}$)}\\ 
        &= \frac{2x}{x^2-y^2} - \frac{2x}{x^2+y^2} && \text{(collect like terms; $(a+b)(a-b) = a^2-b^2$)}\\ 
        &= \frac{2x(x^2+y^2)-2x(x^2-y^2)}{(x^2-y^2)(x^2+y^2)} && \text{($\frac{a}{b} - \frac{p}{q} = \frac{aq-pb}{bq}$)}\\ 
        &= \frac{2x^3+2xy^2-2x^3+2xy^2}{(x^2-y^2)(x^2+y^2)} && \text{(expand numerator)}\\ 
        &= \frac{4xy^2}{x^4+x^2y^2-y^2x^2-y^4} && \text{(collect like terms; expand denominator)}\\ 
        &= \frac{4xy^2}{x^4-y^4} && \text{(collect like terms)}
    \end{align*}

    \problem Simplify---
    \[
        \left(\frac{x}{y} + \frac{y}{x}\right)\left(\frac{a}{b} + \frac{b}{a}\right) - \left(\frac{x}{y} - \frac{y}{x}\right)\left(\frac{a}{b} - \frac{b}{a}\right).
    \]
    \answer
    \begin{align*}
        \left(\frac{x}{y} + \frac{y}{x}\right)\left(\frac{a}{b} + \frac{b}{a}\right) - \left(\frac{x}{y} - \frac{y}{x}\right)\left(\frac{a}{b} - \frac{b}{a}\right) &\\ 
        = \frac{x^2+y^2}{xy} \times \frac{a^2+b^2}{ab} - \frac{x^2-y^2}{xy} \times \frac{a^2-b^2}{ab} && \text{($\frac{a}{b} \pm \frac{p}{q} = \frac{aq \pm bp}{bq}$)}\\ 
        = \frac{(x^2+y^2)(a^2+b^2)}{abxy} - \frac{(x^2-y^2)(a^2-b^2)}{abxy} && \text{($\frac{a}{b} \times \frac{p}{q} = \frac{ap}{bq}$)}\\ 
        = \frac{(x^2+y^2)(a^2+b^2) - (x^2-y^2)(a^2-b^2)}{abxy} && \text{($\frac{a}{b} - \frac{c}{b} = \frac{a-c}{b}$)}\\ 
        = \frac{a^2x^2 + b^2x^2 + a^2y^2 + b^2y^2 - (a^2x^2 - b^2x^2 - a^2y^2 + b^2y^2)}{abxy} && \text{(expand numerator)}\\ 
        = \frac{a^2x^2 + b^2x^2 + a^2y^2 + b^2y^2 - a^2x^2 + b^2x^2 + a^2y^2 - b^2y^2}{abxy} && \text{(remove brackets)}\\ 
        = \frac{2b^2x^2 + 2a^2y^2}{abxy} && \text{(collect like terms)}\\ 
        = \frac{2(b^2x^2 + a^2y^2)}{abxy} && \text{(extract common factor $2$)}
    \end{align*}

    \problem Simplify---
    \[
        \left(\frac{a}{b} - \frac{b}{a}\right) \div \left(\frac{a^2}{b^2} - \frac{b^2}{a^2}\right).
    \]
    \answer 
    \begin{align*}
        \MoveEqLeft
        \left(\frac{a}{b} - \frac{b}{a}\right) \div \left(\frac{a^2}{b^2} - \frac{b^2}{a^2}\right)\\ 
        &= \frac{a^2-b^2}{ab} \div \frac{a^4-b^4}{a^2b^2} && \text{($\frac{a}{b} - \frac{p}{q} = \frac{aq-pb}{bq}$)}\\ 
        &= \frac{a^2-b^2}{ab} \times \frac{a^2b^2}{a^4-b^4} && \text{($\frac{a}{b} \div \frac{p}{q} = \frac{a}{b} \times \frac{q}{p}$)}\\ 
        &= \frac{a^2-b^2}{ab} \times \frac{a^2b^2}{(a^2-b^2)(a^2+b^2)} && \text{($a^2-b^2 = (a-b)(a+b)$)}\\ 
        &= \frac{a^2b^2(a^2-b^2)}{ab(a^2-b^2)(a^2+b^2)} && \text{($\frac{a}{b} \times \frac{p}{q} = \frac{ap}{bq}$)}\\ 
        &= \frac{ab}{a^2+b^2} && \text{($\frac{ac}{bc} = \frac{a}{b}$)}
    \end{align*}

    \problem Simplify---
    \[
        \frac{a(a-b) - b(a+b)}{\dfrac{a}{a+b} - \dfrac{b}{a-b}}.
    \]
    \answer 
    \begin{align*}
        \MoveEqLeft
        \frac{a(a-b) - b(a+b)}{\dfrac{a}{a+b} - \dfrac{b}{a-b}}\\ 
        &= \big(a(a-b) - b(a+b)\big) \div \left(\frac{a}{a+b} - \frac{b}{a-b}\right) && \text{($\frac{a}{b} = a \div b$)}\\ 
        &= \big(a(a-b) - b(a+b)\big) \div \frac{a(a-b) - b(a+b)}{(a+b)(a-b)} && \text{($\frac{a}{b} - \frac{p}{q} = \frac{aq-pb}{bq}$)}\\ 
        &= \big(a(a-b) - b(a+b)\big) \times \frac{(a+b)(a-b)}{a(a-b) - b(a+b)} && \text{($a \div \frac{p}{q} = a \times \frac{q}{p}$)}\\ 
        &= (a+b)(a-b) && \text{($\frac{ac}{bc} = \frac{a}{b}$)}\\ 
        &= a^2-b^2 && \text{($(a+b)(a-b) = a^2-b^2$)}
    \end{align*}

    \problem Simplify---
    \[
        \frac{1-x}{1+x+x^2} - \frac{1+x}{1-x+x^2}.
    \]

    \answer 
    \begin{align*}
        \MoveEqLeft
        \frac{1-x}{1+x+x^2} - \frac{1+x}{1-x+x^2}\\ 
        &= \frac{(1-x)(1-x+x^2) - (1+x)(1+x+x^2)}{(1+x+x^2)(1-x+x^2)} && \text{($\frac{a}{b} - \frac{p}{q} = \frac{aq-pb}{bq}$)}\\ 
        &= \frac{1-2x+2x^2-x^3 - (1+2x+2x^2+x^3)}{(1+x+x^2)(1-x+x^2)} && \text{(expand numerator)}\\ 
        &= \frac{-4x-2x^3}{(1+x^2)^2-x^2} && \text{($(p+q)(p-q) = p^2-q^2$)}\\ 
        &= \frac{-4x-2x^3}{1+x^2+x^4} && \text{(expand denominator)}
    \end{align*}

    \problem Simplify---
    \[
        \frac{1 + \dfrac{x}{1+x}}{x + \dfrac{1}{1+x}} \div \frac{(x+1)^2 - x^2}{x^2+x+1}.
    \]
    \answer 
    \begin{align*}
        \MoveEqLeft
        \frac{1 + \dfrac{x}{1+x}}{x + \dfrac{1}{1+x}} \div \frac{(x+1)^2 - x^2}{x^2+x+1}\\ 
        &= \frac{\dfrac{1+2x}{1+x}}{\dfrac{x^2+x+1}{1+x}} \div \frac{2x+1}{x^2+x+1} && \text{(common denominator $1+x$; expand)}\\ 
        &= \frac{1+2x}{1+x} \times \frac{1+x}{x^2+x+1} \div \frac{2x+1}{x^2+x+1} && \text{($\frac{a}{b} \div \frac{p}{q} = \frac{a}{b} \times \frac{q}{p}$)}\\ 
        &= \frac{1+2x}{x^2+x+1} \times \frac{x^2+x+1}{2x+1} && \text{($\frac{ac}{bc} = \frac{a}{b}$)}\\ 
        &= 1 && \text{($\frac{a}{a} = 1$)}
    \end{align*}

    \problem Simplify---
    \[
        \frac{a^2+b^2}{(a+b)^2} + \frac{\dfrac{2}{ab}\left(\dfrac{1}{a} + \dfrac{1}{b}\right)}{\left(\dfrac{1}{a} + \dfrac{1}{b}\right)^{3}}.
    \]
    \answer 
    \begin{align*}
        \MoveEqLeft
        \frac{a^2+b^2}{(a+b)^2} + \frac{\dfrac{2}{ab}\left(\dfrac{1}{a} + \dfrac{1}{b}\right)}{\left(\dfrac{1}{a} + \dfrac{1}{b}\right)^{3}}\\ 
        &= \frac{a^2+b^2}{(a+b)^2} + \frac{2}{ab} \times \frac{1}{\left(\dfrac{1}{a} + \dfrac{1}{b}\right)^{2}} && \text{($\frac{ac}{bc} = \frac{a}{b}$, cancelling $\tfrac{1}{a}+\tfrac{1}{b}$)}\\ 
        &= \frac{a^2+b^2}{(a+b)^2} + \frac{2}{ab} \times \frac{1}{\left(\dfrac{a+b}{ab}\right)^{2}} && \text{($\frac{1}{a} + \frac{1}{b} = \frac{a+b}{ab}$)}\\ 
        &= \frac{a^2+b^2}{(a+b)^2} + \frac{2}{ab} \times \frac{a^2b^2}{(a+b)^2} && \text{($\left(\frac{a}{b}\right)^m = \frac{a^m}{b^m}$)}\\ 
        &= \frac{a^2+b^2}{(a+b)^2} + \frac{2ab}{(a+b)^2} && \text{($\frac{ac}{bc} = \frac{a}{b}$)}\\ 
        &= \frac{a^2+2ab+b^2}{(a+b)^2} && \text{($\frac{a}{b} + \frac{c}{b} = \frac{a+c}{b}$)}\\ 
        &= \frac{(a+b)^2}{(a+b)^2} && \text{($a^2+2ab+b^2 = (a+b)^2$)}\\ 
        &= 1 && \text{($\frac{a}{a} = 1$)}
    \end{align*}

    \problem Show that---
    \[
        \frac{x^4}{a^2b^2} + \frac{(x^2-a^2)^2}{a^2(a^2-b^2)} - \frac{(x^2-b^2)^2}{b^2(a^2-b^2)}
    \]
    is independent of $x$.
    \begin{proof}
        Taking $a^2b^2(a^2-b^2)$ as a common denominator,
        \[
            \frac{x^4}{a^2b^2} + \frac{(x^2-a^2)^2}{a^2(a^2-b^2)} - \frac{(x^2-b^2)^2}{b^2(a^2-b^2)}
                = \frac{\mathrm{N}}{a^2b^2(a^2-b^2)},
        \]
        where
        \begin{align*}
            \mathrm{N} &= x^4(a^2-b^2) + b^2(x^2-a^2)^2 - a^2(x^2-b^2)^2\\ 
                       &= x^4(a^2-b^2) + b^2(x^4-2a^2x^2+a^4)\\ 
                       &\qquad {}- a^2(x^4-2b^2x^2+b^4) && \text{($(p-q)^2 = p^2-2pq+q^2$)}\\ 
                       &= a^2x^4-b^2x^4 + b^2x^4-2a^2b^2x^2+a^4b^2\\ 
                       &\qquad {}- a^2x^4+2a^2b^2x^2-a^2b^4 && \text{(remove brackets)}\\ 
                       &= a^4b^2-a^2b^4 && \text{(collect like terms)}\\ 
                       &= a^2b^2(a^2-b^2) && \text{(extract common factor $a^2b^2$)}
        \end{align*}
        Then the expression is
        \[
            \frac{a^2b^2(a^2-b^2)}{a^2b^2(a^2-b^2)} = 1.
        \]
        From this we see that all terms of $x$ cancel, so its value is always 1, thus the expression is independent of $x$.
    \end{proof}

    \problem Simplify---
    \[
        \frac{a}{b - \dfrac{c}{d - \dfrac{e}{f}}}.
    \]
    \answer 
    \begin{align*}
        \MoveEqLeft
        \frac{a}{b-\dfrac{c}{d-\dfrac{e}{f}}}\\ 
        &= a \div \left[b - \left(c \div \left(d - \dfrac{e}{f}\right)\right)\right] && \text{($\frac{a}{b} = a \div b$)}\\ 
        &= a \div \left[b - \left(c \div \dfrac{df-e}{f}\right)\right] && \text{($a - \frac{p}{q} = \frac{aq-p}{q}$)}\\ 
        &= a \div \left[b - \dfrac{cf}{df-e}\right] && \text{($a \div \frac{p}{q} = \frac{aq}{p}$)}\\ 
        &= a \div \left[\dfrac{b(df-e)-cf}{df-e}\right] && \text{($a - \frac{p}{q} = \frac{aq-p}{q}$)}\\ 
        &= \frac{a(df-e)}{b(df-e)-cf} && \text{($a \div \frac{p}{q} = \frac{aq}{p}$)}\\ 
        &= \frac{adf-ae}{bdf-be-cf} && \text{(expand)}
    \end{align*}

    \problem Simplify---
    \[
        \frac{1}{a - 2b - \dfrac{1}{a - 2b - \dfrac{1}{a-2b}}}.
    \]
    \answer 
    Writing $t = a-2b$,
    \begin{align*}
        \MoveEqLeft
        \frac{1}{a - 2b - \dfrac{1}{a - 2b - \dfrac{1}{a-2b}}}\\ 
        &= \frac{1}{t - \dfrac{1}{t - \dfrac{1}{t}}} && \text{(replace variables)}\\ 
        &= \frac{1}{t - \dfrac{1}{\dfrac{t^2-1}{t}}} && \text{($\frac{a}{b} - \frac{p}{q} = \frac{aq-pb}{bq}$)}\\ 
        &= \frac{1}{t - \dfrac{t}{t^2-1}} && \text{($a \div \frac{p}{q} = \frac{aq}{p}$)}\\ 
        &= \frac{1}{\dfrac{t(t^2-1)-t}{t^2-1}} && \text{($\frac{a}{b} - \frac{p}{q} = \frac{aq-pb}{bq}$)}\\ 
        &= \frac{t^2-1}{t^3-2t} && \text{(expand; $a \div \frac{p}{q} = \frac{aq}{p}$)}\\ 
        &= \frac{(a-2b)^2-1}{(a-2b)^3-2(a-2b)} && \text{(restore variables)}
    \end{align*}

    \problem Simplify---
    \[
        \frac{a+b}{a + b + \dfrac{1}{a - b + \dfrac{1}{a+b}}}.
    \]
    \answer 
    \begin{align*}
        \MoveEqLeft
        \frac{a+b}{a + b + \dfrac{1}{a - b + \dfrac{1}{a+b}}}\\ 
        &= \frac{a+b}{a + b + \dfrac{1}{\dfrac{(a-b)(a+b)+1}{a+b}}} && \text{($\frac{a}{b} + \frac{p}{q} = \frac{aq+pb}{bq}$)}\\ 
        &= \frac{a+b}{a + b + \dfrac{a+b}{a^2-b^2+1}} && \text{($(a+b)(a-b) = a^2-b^2$; $a \div \frac{p}{q} = \frac{aq}{p}$)}\\ 
        &= \frac{a+b}{\dfrac{(a+b)(a^2-b^2+1)+(a+b)}{a^2-b^2+1}} && \text{($\frac{a}{b} + \frac{p}{q} = \frac{aq+pb}{bq}$)}\\ 
        &= \frac{a+b}{\dfrac{(a+b)(a^2-b^2+2)}{a^2-b^2+1}} && \text{(extract common factor $a+b$)}\\ 
        &= \frac{(a+b)(a^2-b^2+1)}{(a+b)(a^2-b^2+2)} && \text{($a \div \frac{p}{q} = \frac{aq}{p}$)}\\ 
        &= \frac{a^2-b^2+1}{a^2-b^2+2} && \text{($\frac{ac}{bc} = \frac{a}{b}$)}
    \end{align*}
\end{problems}

\subsection{Exercises IV}

\begin{problems}[from=1,step=1]
    \problem If the two fractions $\mathrm{A}/\mathrm{B}$, $a/b$ be equal, and the latter be at its lowest terms, prove that $\mathrm{A} = \mu a$, $\mathrm{B} = \mu b$, where $\mu$ is an integer.

    \begin{proof}
        Suppose $\mathrm{A}$, $\mathrm{B}$, $a$, $b$ integers with $\mathrm{A}/\mathrm{B} = a/b$, and $a/b$ at its lowest terms; that is, $a$ and $b$ have no common measure save unity, so that $a$ is prime to $b$ (\S 6).
        Multiplying across by $b\mathrm{B}$,
        \[
            \mathrm{A}b = a\mathrm{B};
        \]
        hence $a$ divides $\mathrm{A}b$, and being prime to $b$ it must divide $\mathrm{A}$ (\S 10, Cor.\ 1).
        Write $\mathrm{A} = \mu a$, $\mu$ an integer.
        Substituting, $\mu a b = a\mathrm{B}$, and cancelling $a$ (which is not zero), $\mathrm{B} = \mu b$---the same $\mu$ serving in both.
    \end{proof}

    \problem Prove that the sum or difference of two odd numbers is always even; the sum or difference of an odd and an even number always odd; the product of any number of odd numbers always odd; the quotient of one odd number by another always odd, if it be integral.

    \begin{proof}
        Let $a,b,c$ be integers where $a$ is even and $b$ and $c$ are odd.
        Then $a = 2m$, $b = 2n + 1$, and $c = 2l + 1$ for some integers $m,n,l$ (Cor., \S 12).
        Next, we prove the propositions individually.

        \begin{enumerate}
            \item \textit{Sum or difference of two odd numbers is always even}. Take $b + c$, the sum of two odd numbers, this can be written $b + c = (2n+1)+(2l+1) = 2n + 1 + 2l + 1 = 2n + 2l + 2 = 2(n + l + 1)$. Thus, the sum is always even. Now take the difference of two odd numbers, in either order: $b - c = (2n+1)-(2l+1) = 2n + 1 - 2l - 1 = 2n - 2l = 2(n-l)$, and $c - b = (2l+1)-(2n+1) = 2l + 1 - 2n - 1 = 2l - 2n = 2(l-n)$. Thus the difference is always even.

            \item \textit{Sum or difference of an odd and an even number is always odd}. Take $a + b$, the sum of an even number and an odd number, this can be written $a + b = 2m + (2n+1) = 2(m+n) + 1$. Thus, the sum is always odd. Now take the difference, in either order: $a - b = 2m - (2n+1) = 2(m-n-1) + 1$, and $b - a = (2n+1) - 2m = 2(n-m) + 1$. Thus the difference is always odd.

            \item \textit{The product of any number of odd numbers is always odd}. First take two of them, $bc = (2n+1)(2l+1) = 4nl + 2n + 2l + 1 = 2(2nl + n + l) + 1$, which is odd. Now suppose the product of $k$ odd numbers to be odd, say $2p + 1$, and let $2q + 1$ be a further odd number. Then $(2p+1)(2q+1) = 2(2pq + p + q) + 1$, which is odd again. Hence the product of $k+1$ odd numbers is odd whenever the product of $k$ of them is; and since the product of two is odd, the product of any number of odd numbers is odd.

            \item \textit{The quotient of one odd number by another is always odd, if it be integral}. Let $b/c$ be integral, say $b = cq$, where $q$ is an integer. Were $q$ even, $q = 2k$ say, we should have $b = 2ck$, that is, $b$ even, contrary to hypothesis. Hence $q$ is odd.
        \end{enumerate}
    \end{proof}


    \problem If $a$ be prime to $b$, then---
    \begin{parts}
        \item $(a+b)^m$ and $(a-b)^m$ have at most the G.C.M. $2^m$;
        \item $a^m + b^m$ and $a^m - b^m$ have at most the G.C.M. 2;
        \item $a+b$ and $a^2 + b^2 - ab$ have at most the G.C.M. 3.
    \end{parts}

    \begin{lemma*}[A]
        If $x$ be prime to $y$, and $g$ divide $x+y$, then $g$ is prime to $x$ and to $y$.
    \end{lemma*}
    \begin{proof}
        Let $e$ be a common measure of $g$ and $x$.
        Then $e$ divides $x+y$ and $x$. 
        Therefore $e$ divides $(x+y) - x = y$ (\S 8).
        Thus  $e$ is a common measure of $x$ and $y$.
        Because $x$ and $y$ are prime to each other, it follows $e$ is unity.
        Hence $g$ is prime to $x$, and in like manner $g$ is prime to $y$.
    \end{proof}

    \begin{lemma*}[B]
        If $x$ be prime to $y$, the G.C.M. of $x+y$ and $x-y$ divides $2$.
    \end{lemma*}
    \begin{proof}
        Let $g$ be the G.C.M. of $x+y$ and $x-y$.
        Then $g$ divides $(x+y) + (x-y) = 2x$ (\S 8).
        Since $g$ is prime to $x$ (Lemma A), it follows $g$ divides 2 (\S 10, Cor.1).
    \end{proof}

    \begin{proof}[Proof of (a)]
        Let $g$ be the G.C.M. of $a+b$ and $a-b$.
        Then $g$ divides 2 (Lemma B).
        Next, write $a+b = gp$ and $a-b = gq$ for integers $p,q$.
        Then $p$ is prime to $q$; for if they had a common factor $e$ greater than unity, $ge$ would be a common factor of $a+b$ and $a-b$ greater than their G.C.M.
        Hence $p^m$ is prime to $q^m$ (\S 10, Cor. 4), and integers $\mathrm{A}$, $\mathrm{B}$ can be found such that $\mathrm{A}p^m - \mathrm{B}q^m = \pm 1$ (\S 14, Cor.\ 2).
        Thus $\mathrm{A}g^mp^m - \mathrm{B}g^mq^m = \pm g^m$.
        Now $(a+b)^m = g^mp^m$ and $(a-b)^m = g^mq^m$, so that every common factor of $(a+b)^m$ and $(a-b)^m$ divides this algebraic sum of their multiples (\S 8), that is, divides $g^m$.
        But $g$ divides $2$, and therefore $g^m$ divides $2^m$; hence the G.C.M. of $(a+b)^m$ and $(a-b)^m$ is at most $2^m$.
    \end{proof}

    \begin{proof}[Proof of (b)]
        Since $a$ is prime to $b$, $a^m$ is prime to $b^m$ (\S 10, Cor. 4).
        Hence, by Lemma B applied to $a^m$ and $b^m$, the G.C.M. of $a^m + b^m$ and $a^m - b^m$ divides $2$, that is, is at most $2$.
    \end{proof}

    \begin{proof}[Proof of (c)]
        Let $g$ be the G.C.M. of $a+b$ and $a^2 + b^2 - ab$.
        Since $g$ divides $a+b$ it divides $(a+b)^2$, and therefore divides $(a+b)^2 - (a^2 + b^2 - ab) = 3ab$ (\S 8).
        Since $g$ divides $a+b$, it follows $g$ is prime to $a$ and $b$ (Lemma A).
        Hence $g$ is prime to their product $ab$ (\S 10, Cor. 2).
        But $g$ divides $3ab$ and is prime to $ab$, so $g$ divides $3$ (\S 10, Cor. 1), that is, the G.C.M. of $a+b$ and $a^2 + b^2 - ab$ is at most $3$.
    \end{proof}

    \problem The difference of the squares of any two odd numbers is exactly divisible by 8.

    \problem The sum of the squares of three consecutive odd numbers increased by 1 is a multiple of 12.

    \problem If each of two fractions be at its lowest terms, neither their sum nor their difference can be an integer unless the denominators be equal.

    \problem Resolve 45738 and 297675 into their prime factors.

    \problem Find the G.C.M. of 54643 and 91319, using negative remainders whenever it is of advantage to do so.

    \problem Prove that the L.C.M. of two integers is the quotient of their product by the G.C.M.

    \problem If $g_1$, $g_2$, $g_3$ be the G.C.M.'s, $l_1$, $l_2$, $l_3$ the L.C.M.'s, of $b$ and $c$, $c$ and $a$, $a$ and $b$ respectively, $\mathrm{G}$ the G.C.M., and $\mathrm{L}$ the L.C.M., of the three $a$, $b$, $c$, show that
    \begin{parts}
        \item $\mathrm{L} = \dfrac{abc\mathrm{G}}{g_1 g_2 g_3}$;
        \item $\dfrac{\mathrm{L}}{\mathrm{G}} = \sqrt{\dfrac{l_1 l_2 l_3}{g_1 g_2 g_3}}$.
    \end{parts}

    \problem When $x$ is divided by $y$, the quotient is $u$ and the remainder $v$; show that, when $x$ and $uy$ are divided by $v$, the remainders are the same, and the quotients differ by unity.
\end{problems}

\end{document}
